\section{Introduction}

Streamflow prediction plays a central role in water resource management, flood mitigation, and watershed operation regulation \cite{GAO2026103148}. In the context of increasing hydrological and climatic uncertainties \cite{immerzeel2010climate}, it is a key support for high-stakes decision-making in areas like optimal water allocation and precise flood risk management \cite{milly2020colorado}. Therefore, the reliability and accuracy of the model are of equal research importance. In recent years, machine learning and deep learning models have demonstrated excellent performance in streamflow prediction \cite{apak2026multi,ZIPPER2026134909}, but the inherent "black-box" nature of these models makes their internal decision-making mechanisms difficult to interpret \cite{Núñez2023_XAI}, significantly limiting the transparency and reliability of the predictions \cite{WangShap2024}, and potentially misguiding actual decision-making. Thus, improving model interpretability while maintaining predictive performance has become a core challenge in hydrological modeling \cite{MardShap2022}, garnering widespread attention and research from the academic community.

To address the "black-box" characteristic of data-driven models, some scholars have attempted to integrate physical knowledge of hydrological processes into neural network architectures to enhance the physical consistency of the models \cite{SCHMID2026124819}. Among these, Physics-Informed Neural Networks (PINN) are one of the most representative approaches \cite{FANG2025_PINN}, with the core idea of embedding hydrological control equations into the model's loss function as constraints, so that the model strictly adheres to physical laws while fitting observed data \cite{SHI2025_PINN,LIU2026_PINN}, thereby improving the physical plausibility of the model outputs. However, physical consistency does not fundamentally resolve the issue of model interpretability \cite{Roscher2020_Explainable}. The model's utilization mechanisms for input variables \cite{NASER2026_Inrter}, as well as the core logic it uses to represent hydrological processes internally, still remain largely ambiguous \cite{Krishnapriyan2021_PINN}, making it difficult to meet the deeper interpretability requirements for practical decision-making.

Another mainstream approach relies on post hoc analysis techniques to explore model interpretability. Among them, Local Interpretable Model-agnostic Explanations (LIME) \cite{Ribeiro2016_LIME_Origin} and SHapley Additive Explanations (SHAP) \cite{lundberg2017_Shap_Origin} are the two most widely used methods. These techniques build local surrogate models to approximate predictions or decompose them based on cooperative game theory \cite{Clark2025_Shap}, quantifying the contribution of each input feature to the model's output. In this way, they partially break the "black-box" barrier of model transparency \cite{Wagan2026_LIME}, providing intuitive tools for analyzing the impact mechanisms of key hydrological variables such as precipitation and evapotranspiration. However, these post hoc explanation methods have inherent limitations: the explanation results heavily depend on local approximations or statistical decomposition processes, which are easily influenced by data distribution characteristics and parameter settings. Furthermore, these methods lack deep integration with hydrological processes, making it difficult to fully reflect the intrinsic logic behind the model's decision-making, and thus the reliability of the explanation results remains to be improved.

Overall, both PINN methods and post hoc explanation techniques have made significant progress in enhancing the interpretability of streamflow prediction models, offering feasible approaches to address the "black-box" problem. However, both methods still have some shortcomings that can be further improved, which lead to concerns when applying interpretability results to guide high-risk watershed management decisions. Specifically, these concerns manifest in three main areas:

First, the interpretability results lack clear verifiability \cite{NEURIPS2021_shap_residuals}. Unlike model predictions, which can be directly compared and validated against observed streamflow data, feature attribution-based interpretability results do not have objective ground truth values, making it difficult to quantitatively verify them through independent experiments \cite{TAKEFUJI2026_Shap_problem}. Existing studies have shown that even if a model's prediction accuracy is high, its interpretability results may only reflect complex statistical fitting relationships or spurious correlations \cite{e24050687_Shap_problem}. Relying solely on feature attribution analysis makes it challenging to determine whether the model has genuinely learned physically meaningful hydrological relationships \cite{Benchmarkingofa_PhysicallyBasedHydrologicModel}, which severely limits the trustworthiness of interpretability results in high-risk decision-making environments.

Second, there is a lack of systematic evaluation of the model's behavior. Most current interpretability studies remain confined to feature-level attribution analysis based on SHAP or LIME \cite{ZHOU2026103254,XIA2026103232}, without thoroughly testing whether the model's response to changes in hydrological conditions aligns with physical logic \cite{10.2166/wcc.2017.149}. In practice, a physically reliable streamflow prediction model should generate consistent physical responses when key driving factors are disturbed \cite{hess-30-1077-2026}—for example, streamflow should steadily increase when precipitation significantly rises \cite{nearing2021role}. Without such perturbation testing or counterfactual validation \cite{simic2025perturbation,ZHANG2025123192}, interpretability results cannot be effectively linked to actual hydrological behavior \cite{read2019process}, and thus fail to provide genuinely valuable insights for decision-making.

Third, there is a lack of systematic evaluation of the transferability of interpretability results. Current interpretability analyses are mostly confined to a single model architecture and rarely focus on whether the extracted hydrological relationships remain valid across different model families \cite{pmlr-v235-wang24j}. If identified feature relationships cannot be reproduced or applied in other models, these relationships are more likely to reflect the parameterization characteristics of a specific model, rather than generalizable hydrological scientific knowledge. Therefore, cross-model validation is crucial for establishing the scientific value of interpretability results \cite{li2022transferability}. It is worth noting that in other fields, systematic interpretability research paradigms have gradually advanced, and standardized paradigms can significantly enhance trust in model results. However, such attempts remain relatively limited in the hydrological modeling domain, and a mature interpretability verification system has yet to be formed.

This study systematically uncovers the internal information flow mechanisms of deep learning-based streamflow prediction models through multidimensional verification methods, confirming the presence of physical knowledge in the model and clarifying that it is the core foundation for achieving accurate streamflow predictions. Based on this, the study focuses on the following key scientific questions: 

\begin{enumerate}
    \item What is the intrinsic relationship between model structure and prediction performance? 
    \item Does physical knowledge exist in deep learning-based hydrological models, and what patterns govern its spatial and feature distributions? 
    \item Can the model generate feature representations that adhere to physical laws for different hydrological scenarios?
    \item Does the physical knowledge learned by the model possess transferability across scenarios and model structures?
\end{enumerate}

To address these scientific questions, this paper proposes an interpretable hydrological model (STEH), aiming to achieve full-chain interpretability from input features, internal model structures, to output results. The main contribution of this study lies not only in the development of an interpretable deep learning framework, but also the establishment of a reference path for hydrological model interpretability verification. Framed as \textbf{\textit{A Deep Dive into Model Structure}}, this work integrates structural analysis, behavioral verification, and cross-model evaluation to move interpretability beyond descriptive attribution toward rigorous assessment of whether model knowledge is \textbf{\textit{verifiable}}, physically meaningful, and \textbf{\textit{transferable}}. In doing so, it provides preliminary insights into a fundamental question: \textbf{\emph{what kind of hydrological knowledge is learned, and is it verifiable and transferable?}}